\documentclass[a4paper,12pt]{article}
\usepackage[margin=1in]{geometry}
\usepackage{amsmath,amssymb,amsthm}
\usepackage{bm}
\usepackage{enumitem}
\usepackage{listings}
\usepackage{xcolor}
\usepackage{hyperref}
\usepackage{cleveref}
\usepackage{booktabs}
\usepackage{array}
\usepackage{graphicx}
\usepackage{float}
\usepackage{caption}
\usepackage{subcaption}
\usepackage{tikz}
\usepackage{tikz-cd}
\usepackage{authblk}

%--- listings style -------------------------------------------------
\lstset{
    basicstyle=\ttfamily\small,
    keywordstyle=\color{blue},
    commentstyle=\color{gray},
    stringstyle=\color{orange},
    breaklines=true,
    showstringspaces=false,
    numbers=left,
    numberstyle=\tiny\color{gray},
    frame=single,
    captionpos=b,
    language=Python
}
%--------------------------------------------------------------------

\title{Socle Filtration via Littlewood--Richardson Coefficients}
\author{Pablo Santos Guerrero}
\date{\today}

\begin{document}
\maketitle
\begin{abstract}
We give a self- contained, rigorous exposition of a Python program that computes the
\((k\!+\!1)^{\text{st}}\) socle layer of a tensor product of induced modules.
The algorithm relies on combinatorial representation theory:
partitions, skew shapes, Littlewood--Richardson (LR) coefficients, and
Yamanouchi (lattice) words.  
All mathematical notions are introduced together with the corresponding
code fragments; the overall flow of the program is described in detail.
\end{abstract}

\tableofcontents
\newpage

%%%%%%%%%%%%%%%%%%%%%%%%%%%%%%%%%%%%%%%%%%%%%%%%%%%%%%%%%%%%%%%%%%%%%%%%%%%%
\section{Mathematical background}
\label{sec:mathbg}

\subsection{Partitions and degree}
A \emph{partition} of a non-negative integer \(n\) is a weakly decreasing
sequence of positive integers
\[
\lambda = (\lambda_{1},\lambda_{2},\dots,\lambda_{\ell}),\qquad
\lambda_{1}\ge \lambda_{2}\ge\cdots\ge \lambda_{\ell}>0,
\]
with \(|\lambda|:=\sum_{i=1}^{\ell}\lambda_{i}=n\).  
In the code this is implemented by

\begin{lstlisting}[caption={Degree of a partition}]
def degree(partition):
    return sum(partition)
\end{lstlisting}

\subsection{Skew shapes}
If \(\mu\subseteq\lambda\) (i.e. \(\mu_{i}\le\lambda_{i}\) for every row),
the \emph{skew shape} \(\lambda/\mu\) is obtained by deleting the Young diagram
of \(\mu\) from that of \(\lambda\). Its row lengths are
\[
(\lambda_{1}-\mu_{1},\;\lambda_{2}-\mu_{2},\dots).
\]
The routine

\begin{lstlisting}[caption={Skew shape as a list of row lengths}]
def subtract_partitions(lam, mu):
    # returns None if mu is not contained in lam
    if len(mu) > len(lam) or any(mu[i] > lam[i] for i in range(len(mu))):
        return None
    return [lam[i] - (mu[i] if i < len(mu) else 0) for i in range(len(lam))]
\end{lstlisting}
produces precisely this list.

\subsection{Littlewood--Richardson coefficients}
Given partitions \(\mu,\nu,\lambda\),
the Littlewood--Richardson coefficient
\(c^{\lambda}_{\mu,\nu}\in\mathbb{Z}_{\ge 0}\) counts the number of
semistandard Young tableaux of shape \(\lambda/\mu\) having content \(\nu\)
and whose reading word is a \emph{Yamanouchi (lattice)} word.
Two variants are required:

\begin{itemize}[nosep]
    \item \textbf{2-order LR coefficient} \(c^{\lambda}_{\mu,\nu}\) -  denoted in the code by \texttt{lrcoefP}.
    \item \textbf{4-order LR coefficient}
    \[
        c^{\lambda}_{\mu_{1},\mu_{2},\mu_{3},\mu_{4}}
        \;=\;
        \sum_{a\vdash d_{12}}\;\sum_{b\vdash d_{34}}
        c^{a}_{\mu_{1},\mu_{2}}\,
        c^{b}_{\mu_{3},\mu_{4}}\,
        c^{\lambda}_{a,b},
    \]
    where \(d_{12}=|\mu_{1}|+|\mu_{2}|\) and \(d_{34}=|\mu_{3}|+|\mu_{4}|\).
    This is implemented by \texttt{lrcoef4}.
\end{itemize}

\subsection{Yamanouchi (lattice) words}
A word \(w=w_{1}w_{2}\dots w_{m}\) over \(\mathbb{N}_{>0}\) is Yamanouchi iff
for every prefix \(w_{1}\dots w_{k}\) and each integer \(i\ge 1\),

\[
\#\{j\le k\mid w_{j}=i\}\ \ge\ 
\#\{j\le k\mid w_{j}=i+1\}.
\]

The test used in the program is

\begin{lstlisting}[caption={Yamanouchi test}]
def is_yamanouchi(word):
    counts = {}
    for x in word:
        counts[x] = counts.get(x, 0) + 1
        for y in range(1, x):
            if counts.get(y, 0) < counts[x]:
                return False
    return True
\end{lstlisting}

---------------------------------------------------------------------

\section{Algorithmic components}
\label{sec:algo}

\subsection{Generating and filtering partitions}
\begin{description}[font=\normalfont\bfseries]
    \item[\texttt{generate\_partitions(n, max\_part=None)}]  
    Recursively enumerates all integer partitions of \(n\).  

    \item[\texttt{subset\_partitions(partition)}]  
    Returns every partition whose degree does not exceed that of the argument;
    it simply calls \texttt{generate\_partitions} for each integer
    \(0\le m\le |\text{partition}|\).

    \item[\texttt{restrict\_partitions(part1\_list, part2\_list, d\_lam)}]  
    From two lists keep only pairs \((a,b)\) satisfying \(|a|+|b|=d_{\lambda}\).

    \item[\texttt{ones\_partition(x)}]  
    Produces the all-ones partition \((\underbrace{1,\dots,1}_{x\text{ times}})\); it
    will later represent the quantities \(p\) and \(q\) appearing in the socle
    formula.
\end{description}

All of these functions are purely combinatorial and have no hidden algebraic
dependencies.

\subsection{2-order Littlewood--Richardson coefficient (\texttt{lrcoefP})}
\begin{lstlisting}[caption={2-order LR coefficient}]
def lrcoefP(mu, nu, lam):
    # 1. Skew shape lam / mu
    skew = subtract_partitions(lam, mu)
    if skew is None:          # mu not contained in lambda
        return 0
    # 2. Degree compatibility
    if sum(skew) != sum(nu):
        return 0
    # 3. Trivial corner cases
    if mu == () and degree(nu) == degree(lam) and nu != lam:
        return 0
    if nu == () and degree(mu) == degree(lam) and mu != lam:
        return 0
    # 4. Enumerate all words of content nu
    entries = weight_to_list(nu)          # e.g. nu =(2,1) -> [1,1,2]
    c = 0
    for perm in set(permutations(entries)):
        if is_yamanouchi(perm):
            c += 1
    return c
\end{lstlisting}
The function follows the definition verbatim:
* it builds the skew shape,
* it checks that the number of boxes matches the size of \(\nu\),
* it enumerates all distinct permutations of the multiset \(\nu\),
* and it counts those that satisfy the Yamanouchi condition.
Because of the exhaustive permutation step the routine is feasible only for
very small partitions (typically \(|\nu|\le 5\)).

\subsection{4-order Littlewood--Richardson coefficient (\texttt{lrcoef4})}
\begin{lstlisting}[caption={4-order LR coefficient}]
def lrcoef4(mu1, mu2, mu3, mu4, lam):
    d_mu1 = degree(mu1); d_mu2 = degree(mu2)
    d_mu3 = degree(mu3); d_mu4 = degree(mu4)
    total1 = d_mu1 + d_mu2          # |a|
    total2 = d_mu3 + d_mu4          # |b|
    parts_a = list(generate_partitions(total1))
    parts_b = list(generate_partitions(total2))
    s = 0
    for a in parts_a:
        for b in parts_b:
            N1 = int(lrcoefP(mu1, mu2, a))   # c^{a}_{mu1,mu2}
            if N1 == 0: continue
            N2 = int(lrcoefP(mu3, mu4, b))   # c^{b}_{mu3,mu4}
            if N2 == 0: continue
            N3 = int(lrcoefP(a, b, lam))     # c^{lambda}_{a,b}
            if N3 == 0: continue
            s += N1 * N2 * N3
    return s
\end{lstlisting}
The double loop over all possible intermediate partitions \(a\) and \(b\) implements
the summation formula displayed earlier.  In a production version one would prune
the loops using the result of \texttt{restrict\_partitions} and cache already
computed LR numbers; the present script keeps the implementation straightforward.

\subsection{Solving the degree balance system (\texttt{solveOne})}
The socle filtration imposes linear relations among the degrees of several
auxiliary partitions \(\delta,\gamma,p,q\).  The system is

\[
\begin{aligned}
x_{1}+2x_{2}+x_{3}+x_{4}&=k,\\
x_{1}+x_{2}+x_{3}&=k_{1}-k_{2},\\
x_{1}+x_{2}+x_{4}&=k_{3}-k_{4},
\end{aligned}
\qquad\text{with }x_{i}\ge0,
\]

where  

\[
x_{1}=|\delta|,\;x_{2}=|\gamma|,\;x_{3}=p,\;x_{4}=q,
\]
and \(k_{1}=|\lambda|,\ k_{2}=|\lambda'|,\ k_{3}=|\mu|,\ k_{4}=|\mu'|\).

The function brute-forces the small search space (the indices are bounded by
\(k\)) and returns every admissible 4-tuple as a dictionary.

\subsection{Socle filtration formula (\texttt{CalcSoc})}
Let \(\lambda,\lambda',\mu,\mu'\) be the partitions supplied by the user and
\(k\) the filtration index.  
Define the set \(\mathcal{S}(k)\) of all solutions of the above linear system.
For a solution \((\delta,\gamma,p,q)\) we form the two 4-order LR coefficients

\[
c^{\lambda}_{\lambda',\,\gamma,\,\delta,\,p},
\qquad
c^{\mu}_{\mu',\,\gamma,\,\delta,\,q}.
\]

The multiplicity of the \((k+1)^{\text{st}}\) socle layer is then

\[
\boxed{
\displaystyle
\sum_{(\delta,\gamma,p,q)\in\mathcal{S}(k)}
c^{\lambda}_{\lambda',\gamma,\delta,p}\;
c^{\mu}_{\mu',\gamma,\delta,q}
}.
\tag{1}
\]

The implementation follows (1) verbatim:

\begin{lstlisting}[caption={Socle computation}]
def CalcSoc(k, lam, lamP, mu, muP):
    k1 = degree(lam);  k2 = degree(lamP)
    k3 = degree(mu);   k4 = degree(muP)

    # Corner case: when a prime partition already has full degree
    if (k1 == k2) or (k3 == k4):
        return 1

    solutions = solveOne(k, k1, k2, k3, k4)
    total = 0
    for sol in solutions:
        d_list = list(generate_partitions(sol["deg delta"]))
        g_list = list(generate_partitions(sol["deg gamma"]))
        p = ones_partition(sol["p"])
        q = ones_partition(sol["q"])
        for d in d_list:
            for g in g_list:
                term1 = lrcoef4(lamP, g, d, p, lam)   # c^{lam}_{lam',gamma,delta,p}
                term2 = lrcoef4(muP, g, d, q, mu)    # c^{mu}_{mu',gamma,delta,q}
                total += term1 * term2
    return total
\end{lstlisting}

If either \(|\lambda|=|\lambda'|\) or \(|\mu|=|\mu'|\) the module is already simple,
hence the multiplicity equals \(1\).

\subsection{Driver routine (\texttt{Master})}
The function \texttt{Master} is a thin wrapper that prints a nicely formatted
header, invokes \texttt{CalcSoc}, and displays the result.

\begin{lstlisting}[caption={Driver}]
def Master(k, lam, lamP, mu, muP, ex_string = ''):
    print(f'''|----------- Calculation for {k+1}-th socle filtration ----------|
            lambda = {lam}
            lambda' = {lamP}
            mu = {mu}
            mu' = {muP}
            k = {k}
          ''')
    print(f'>>>>> {k+1}-th Socle Filtration is : {CalcSoc(k, lam, lamP, mu, muP)} <<<<<')
    print('-------------------- Done --------------------')
\end{lstlisting}

A concrete invocation used for testing is

\begin{verbatim}
Master(4, (1,1), (), (1,1), ())
\end{verbatim}
which computes the 5-th socle layer for \(\lambda=\mu=(1,1)\) and
\(\lambda'=\mu'=\varnothing\); the output is \(4\).

--------------------------------------------------------------------

\section{Complexity considerations and practical tips}
\label{sec:complexity}

\begin{itemize}
    \item \textbf{Partition generation.}
          The number of partitions of \(n\) is the partition function \(p(n)\),
          which grows roughly like \(\exp\bigl(\pi\sqrt{2n/3}\bigr)/(4n\sqrt{3})\).
          Hence exhaustive generation is only realistic for \(n\lesssim 8-10\).

    \item \textbf{LR coefficient \texttt{lrcoefP}.}
          The algorithm enumerates all distinct permutations of the multiset
          described by \(\nu\); its complexity is  
          \(\displaystyle \frac{(\sum \nu_i)!}{\prod \nu_i!}\).
          For \(|\nu|\le 5\) this is still manageable; beyond that a
          combinatorial LR implementation (e.g. via the hive model or
          jeu-de-taquin) is required.

    \item \textbf{Four-order coefficient \texttt{lrcoef4}.}
          Without pruning it performs \(|\mathcal{P}(d_{12})|\times |\mathcal{P}(d_{34})|\)
          calls to \texttt{lrcoefP}.  A cheap optimisation is to
          pre-filter the pair \((a,b)\) using \texttt{restrict\_partitions}
          before the inner loop.

    \item \textbf{Memoisation.}
          Many identical triples \((\mu,\nu,\lambda)\) appear repeatedly.
          Storing already computed LR values in a dictionary reduces the overall
          runtime dramatically.

    \item \textbf{Parallelism.}
          The outer loops over \(\delta\), \(\gamma\) are independent and can be
          dispatched to multiple processes (e.g. via ``multiprocessing.Pool``).

    \item \textbf{Typical use-case.}
          In the original research the program is invoked only for
          small-degree partitions (often \(|\lambda|,|\mu|\le 4\)), where the
          naive enumeration is perfectly acceptable.
\end{itemize}

---------------------------------------------------------------------

\section*{References}
\begin{enumerate}
\item W. Fulton, \emph{Young Tableaux}, Cambridge Univ. Press, 1997.
\item R. P. Stanley, \emph{Enumerative Combinatorics, Vol. 2},
      Cambridge Univ. Press, 1999 (Chapter 7 on LR rule).
\item G. James, A. Kerber, \emph{The Representation Theory of the Symmetric Group},
      Addison- Wesley, 1981.
\item J. E. Humphreys, \emph{Representations of Reductive Groups},
      Cambridge Univ. Press, 1991 -  socle filtration background.
\end{enumerate}

\end{document}
